\chapter{交互总线：从内部决策到环境反馈的闭环动力学}
\label{chp:interaction_bus_ch}

\label{chp:teci_cycle}
上一章建立了从离散事件到连续状态、再到离散承诺与慢结构更新的内部循环。本章转向它在环境一侧的延伸：一个智能系统怎样把内部目标变成可执行控制，怎样让控制经过身体、工具与环境产生后果，又怎样把后果重新变成可校验的观测。作者把这条外循环称为 TECI，但不把它规定为所有智能系统都必须采用的唯一模块名。它所指向的是更一般的计算问题：决策变量、执行接口、环境状态和反馈证据属于不同空间，闭环能否成立取决于显式的类型转换、时间尺度、资源约束与因果归因。双纽线、莫比乌斯带、辐射和阻抗在本章只作为组织直觉；论证本身建立在状态空间模型、控制接口、概率核与可检验误差之上。

\section{边界映射：内部目标与外部控制不是同一种对象}
\label{sec:section_6104}

设智能体内部状态为 $x_t\in\mathcal X\subseteq\mathbb R^{n_x}$，任务目标为 $g_t\in\mathcal G$，环境状态为 $s_t\in\mathcal S\subseteq\mathbb R^{n_s}$，可执行控制为 $u_t\in\mathcal U\subseteq\mathbb R^{n_u}$，传感观测为 $y_t\in\mathcal Y\subseteq\mathbb R^{n_y}$。若分量来自不同量纲，先用预先声明的尺度矩阵进行无量纲化；若控制量具有牛顿、伏特或每秒请求数等工程单位，则在接口处保留单位而不能与语义置信度相加。系统在离散决策时刻 $t_k$ 产生承诺 $a_k\in\mathcal A$，执行器在区间 $[t_k,t_{k+1})$ 内依据策略接口
\[
u(t)=\Pi_{\mathrm{act}}\bigl(a_k,x_{t_k},y_{0:k};\theta_{\mathrm{act}}\bigr)
\]
生成控制轨迹。这里 $\Pi_{\mathrm{act}}$ 是工程设计的解码器，不是“意志变成力量”的本体论等式；它必须给出输入类型、输出范围、延迟、饱和和失败码。环境随后服从模型假设
\[
\dot s(t)=f_{\mathrm{env}}\bigl(s(t),u(t),w(t),t\bigr),\qquad
y_k=h_{\mathrm{obs}}\bigl(s(t_k),v_k,t_k\bigr),
\]
其中 $w$ 与 $v$ 分别表示过程扰动和观测噪声。若采用离散实现，应声明步长 $\Delta t_k=t_{k+1}-t_k$，例如 $s_{k+1}=F_{\Delta t_k}(s_k,u_k,w_k)$；当步长、输入或环境结构变化时，这些变化都进入 $F$，不能把固定场景下的差分直接解释为普遍导数。

这组定义揭示了边界的第一条原则：内部目标与外部控制之间通常不存在全局双射，更不存在由语义相似自动推出的拓扑共轭。多个内部状态可以产生同一动作，一个动作也可能因身体参数、工具状态和环境约束而产生不同后果。我们只在任务相关的局部区域内要求接口保持必要关系。令 $q:\mathcal S\to\mathcal Z$ 提取与任务有关的环境变量，$r:\mathcal X\to\widehat{\mathcal Z}$ 给出内部预测；若存在容差 $\varepsilon$，使闭环运行时 $d_{\mathcal Z}(q(s_k),r(x_k))\leq\varepsilon$，最多可以说二者在指定任务、度量和时间窗内实现了可用对齐，不能据此宣称主观世界与客观世界整体同构。原稿用莫比乌斯单侧曲面表达“感知路径连续转入行动路径”，作者保留其闭环直觉，但把可检验内容改写为两个有方向的类型化通道：行动通道 $\mathcal A\to\mathcal U\to\mathcal S$ 与感知通道 $\mathcal S\to\mathcal Y\to\mathcal X$。二者通过时间戳、对象身份、因果标记和任务上下文绑定，而不是通过神秘的本体翻转。

边界的第二条原则是可撤销性。任何把内部承诺输出到环境的接口，都应附带前置条件 $c_k$、作用域 $B_k$、预算 $b_k$、预期后果分布 $P(\hat y_{k+1:k+H}\mid a_k)$ 和补偿操作 $a_k^{-1}$；不可逆操作则必须提高确认阈值并缩小试探幅度。以机器人拿杯子为例，“把杯子移到桌面左侧”是目标描述，关节力矩序列才是控制，抓取成功、杯体破损与人员接近是不同观测事件。内部语言模型即使正确生成目标，也不等于控制器已经满足摩擦锥、碰撞距离与电机限幅。类似地，软件 Agent 的“更新数据库”必须经过权限、事务、幂等键和回滚接口。由此，所谓知行合一不是取消层间差异，而是让每次跨界转换都有类型契约、证据链和责任边界。

边界的第三条原则是有限可识别性。只观察一次成功动作，通常不能确定成功来自策略、偶然扰动还是外部协助。TECI 因此需要在动作前保存预测，在动作后记录实际观测，并维护未观测混杂因素的假设集合。对可控系统可以采用随机试探或对照动作估计局部响应；对高风险系统则使用数字孪生、离线回放或保守集合传播。只有当接口定义、观察窗、干预条件和替代解释都被陈述后，行动—反馈关系才具有工程意义。这个约束也划清 HSF-HD 与 AgentOS 的层级：前者研究一般智能的状态—边界—环境动力学，后者可以用 Boundary、CCM、TCE 和外部记忆 Me 实现其中一种可审计接口，但这些组件并不构成所有智能系统的先验唯一结构。

\section{TECI 四阶段：发起、传播、作用与回执}
\label{sec:teci_radiation_dynamics}

为了保留原稿“辐射—传播—散射—干涉”的机制顺序，同时避免把波动术语当作证明，作者把 TECI 写成四阶段事务。第一阶段是发起：内部循环形成承诺 $a_k$ 后，行动解码器根据当前状态、资源预算和安全约束选择控制包
\[
c_k=(\mathrm{id}_k,t_k,a_k,u_{k:k+H},B_k,b_k,\widehat Y_k,\rho_k),
\]
其中 $\mathrm{id}_k$ 是唯一身份，$B_k$ 是允许影响的边界，$b_k$ 是时间、调用次数或物理能量预算，$\widehat Y_k$ 是预测观测集合，$\rho_k$ 是风险等级。对 AgentOS 实例，$a_k$ 可以来自 $h(t)$ 中经 $\Psi$ 与 SPC 绑定的任务承诺，TCE 决定探索强度，CCM 检查认知与调用负载，Boundary 决定哪些数据和工具可以越界；但一般系统只需实现等价的约束功能。发起阶段失败的典型情形包括目标尚未消歧、执行器能力与动作不匹配、资源不足和权限无效。此时正确行为是拒绝或降级，而不是把高置信语言输出误当成已经发生的环境改变。

第二阶段是传播。控制包离开决策器后，要经过驱动器、网络、工具链或组织流程。可把这段链路表示为带状态的信道
\[
z_{j+1}=C_j(z_j,\eta_j,\delta_j),\qquad j=0,\ldots,m-1,
\]
其中 $z_0=c_k$，$\eta_j$ 表示丢包、量化、权限或人工解释等扰动，$\delta_j$ 表示延迟。传播不是无损的“幺正演化”：重试可能复制操作，缓存可能使用旧参数，执行者可能改变动作语义。因而每一跳应保留身份、版本与截止时间，并把不可接受的偏差显式转成异常。对于机械体，传播链可能是规划器—伺服器—驱动电路—关节；对于软件体，可能是模型—工具路由—API—数据库；对于组织智能，可能是战略—制度—岗位—现场执行。三者的物理实现不同，但都可检验命令是否到达、是否被改写以及延迟是否超过闭环稳定区间。

第三阶段是作用。环境转移核
\[
P(s_{k+1}\mid s_k,u_k,e_k)
\]
描述控制在上下文 $e_k$ 中改变状态的概率。原稿用阻抗匹配说明“同样的意图遇到不同环境会得到不同结果”，这个直觉可重建为可达性与增益问题。定义任务输出 $z_k=q(s_k)$，局部线性化为 $\Delta z_k\approx B_k^{\mathrm{env}}\Delta u_k+d_k$；矩阵 $B_k^{\mathrm{env}}$ 的秩与条件数刻画当前接口能否有效控制任务维度，$d_k$ 汇总未建模扰动。若控制方向落在近似零空间，再增大动作幅度也可能只增加代价而不改善任务；若增益过高，小误差又会被放大。拿杯子的成功依赖抓取位姿、摩擦、杯中液体和他人干预，数据库更新则依赖锁、约束和并发事务。这里的“匹配”不是统一波阻抗公式，而是针对具体系统估计可控方向、灵敏度和饱和边界。

第四阶段是回执。系统从环境取得 $y_{k+1:k+H}$，先验证其时间戳与对象身份，再与发起阶段保存的预测分布比较。可定义标准化残差
\[
e_k=W_k^{1/2}\bigl(\phi(y_{k+1:k+H})-\hat z_k\bigr),
\]
其中 $\phi$ 是从原始观测到任务变量的读出，$W_k$ 是无量纲权重或协方差逆矩阵，$\hat z_k$ 是预测值。$\|e_k\|$ 小只说明在当前度量下观测接近预测，不自动证明任务成功；成功还要检查约束、迟发副作用和目标版本。$\|e_k\|$ 大也不必然意味着策略错误，它可能来自传感器故障、对象绑定错误或环境突变。系统因此把残差分流为状态修正、模型更新、接口故障和目标重审四类。原稿的“回波干涉”在这里成为预测—观测比较，所谓惊奇激波成为带来源的异常事件。只有异常被正确绑定到导致它的动作、对象和时间窗，下一轮内部循环才不会把无关反馈写入 $E$ 或错误固化到 $\Theta$。

四阶段共同构成一项可追踪事务：发起回答“准备做什么”，传播回答“控制怎样到达”，作用回答“环境怎样响应”，回执回答“我们凭什么判断”。在开放环境中，闭环还必须允许超时、取消、部分成功和补偿，而不能假设每次动作都完整返回。雷达隐喻之所以有用，是它提醒系统保存发射模板并等待回波；它的边界也同样清楚：社会规则、数据库事务和工具调用不必服从波动方程，不能从相位语言推导其因果机制。

\section{双循环耦合：从一次事务到持续适应}
\label{sec:section_4121}

TDCI 与 TECI 可以用双纽线表示两个在当前界面相交的循环，但数学对象应是耦合的快慢系统，而不是一条具有神秘翻转性质的曲面。令 $x_k$ 表示认知工作状态，$s_k$ 表示环境状态，$\theta_k$ 表示慢结构参数，$m_k$ 表示情景记录。一个最小离散模型为
\[
\begin{aligned}
x_{k+1}&=F_x(x_k,y_k,g_k,\theta_k),\\
u_k&=\Pi(x_k,g_k;\theta_k),\\
s_{k+1}&=F_s(s_k,u_k,w_k),\\
y_{k+1}&=H(s_{k+1},v_{k+1}),\\
m_{k+1}&=F_m(m_k,x_k,u_k,y_{k+1}),\\
\theta_{k+1}&=\theta_k+\alpha_k G(\theta_k,m_{0:k+1},g_{0:k+1}).
\end{aligned}
\]
各变量的维度与单位由具体实现给出，$\alpha_k$ 是慢更新步长；若目标 $g_k$、接口 $\Pi$ 或环境转移 $F_s$ 随时间变化，它们必须作为时变量进入更新。前三行形成运行态快闭环，后两行形成经验与结构的慢闭环。对三相记忆实例，$x_k$ 可对应 $h(t)$ 的有限工作界面，$m_k$ 对应 $E$ 中带时间与身份的事件，$\theta_k$ 对应 $\Theta$ 的生成约束；Me 则提供外部读取与写入通道。这个映射服务于 AgentOS 工程说明，并不意味着任何智能都必须以相同存储模块实现。

双循环的关键不是“内部熵减等于外部做功”，而是预算与信息在不同账本中闭合。定义任务损失 $L_k$、计算代价 $C_k^{\mathrm{comp}}$、接口代价 $C_k^{\mathrm{act}}$、风险代价 $C_k^{\mathrm{risk}}$ 和可选的焦耳能耗 $E_k^{\mathrm J}$。这些量单位不同，只有经明确的权重或约束才能共同进入优化：
\[
\min_{\pi}\;\mathbb E\!\left[\sum_{k=0}^{T}\gamma^k L_k\right],\quad
\text{s.t. }\sum C_k^{\mathrm{comp}}\leq B_c,\;
\sum C_k^{\mathrm{act}}\leq B_a,\;
P(\text{unsafe})\leq\delta.
\]
焦耳能耗若被测量，可以另设 $\sum E_k^{\mathrm J}\leq B_E$；不能把信息熵、活动强度、经济成本和热力学能量都写成一个“能量”。耗散结构理论提供开放系统必须交换物质、能量和信息的背景，但不会自动给出具体控制律。系统是否维持远离平衡态，要由资源流、故障恢复和状态稳定性证据支持，而非由“呼吸”修辞直接证明。

持续适应还需要区分局部结构网络与全局分布表示。局部网络保存对象、工具、约束和事件之间可追踪的节点—边结构，适合身份一致性与因果审计；全局分布表示编码跨模态相似性、背景语境与柔性泛化，适合从高维输入形成候选。两者通过四类接口联合：任务相关身份把分布候选落到具体对象，绑定记录某次动作涉及哪些对象与时间，耦合门决定哪些局部约束进入全局状态，读出接口把全局状态转换为可检查的控制变量。人脑区域分工可以提供启发，却不能证明工程模块与脑区机制等同。若绑定失败，系统可能把 A 杯子的反馈归给 B 杯子；若只使用局部图，系统又可能难以处理新情境。CCM 的一般职责正是在有限工作界面中协调这两类表示的激活与卸载，而非把全部知识同时展开。

当反馈连续到来时，快闭环的稳定并不保证慢学习正确。若系统以高增益迅速修正每个残差，可能追逐噪声；若结构更新过慢，又会在环境变化后长期使用过期模型。可用双时间尺度条件表达基本要求：快系统在固定 $\theta$ 下应在相关工作区内输入到状态稳定，慢步长 $\alpha_k$ 相对执行周期足够小，同时变化检测器在分布突变时允许重置或切换模型。这些是模型假设和设计条件，不是对所有非线性开放系统的普遍定理。对高风险动作，还应设置人类或上级控制器的仲裁边界。双纽线因此不是一个保证永不破裂的符号，而是一张审计图：每次内循环给外循环什么承诺，外循环给内循环什么证据，哪些变量在快尺度修正，哪些结构在慢尺度更新。

从这个角度看，所谓智能的“呼吸”是可中断、可恢复的连续事务流。没有感知与内部建模，行动会退化为盲目刺激；没有行动与环境验证，内部推演会失去现实约束；没有情景记录和结构学习，相同错误会反复出现；没有边界与卸载，有限工作界面会被未决事务堵塞。系统的自我连续性并非拓扑孤立子自动产生，而是由身份链、记忆写入规则、目标版本和责任边界共同维持。

\section{双向重构：环境改变、模型更新与任务相关共生}
\label{sec:section_7516}

TECI 的长期结果包含两个方向：智能体改变环境，环境反过来改变智能体。为了避免把“几何扭曲”误写成真实时空曲率，作者把内部结构表示为带参数的预测与控制模型 $M_{\theta}$，把外部结构表示为环境配置 $e\in\mathcal E$。行动可以直接改变任务状态 $s$，也可以改变未来交互的配置，例如制造工具、修路、建立数据库索引或制定组织流程。设配置更新为
\[
e_{k+1}=e_k+\Delta_e(e_k,u_k,\zeta_k),
\]
其中 $\Delta_e$ 只是一类工程转移，不具有统一物理单位；机械位移、权限变化和索引建立应分别建模。把崎岖路径改造成道路的效果，不是抹平黎曼曲率，而是降低后续任务的控制代价、失败率或搜索深度。可用反事实差值
\[
V^{\pi}(s,e')-V^{\pi}(s,e)
\]
评估某次外部重构是否提高了给定策略的长期价值。只有在目标、策略、比较基线和副作用一致时，才能把这个差值归因于环境工程。

反方向的模型更新始于可归因残差。令预测模型给出 $p_{\theta_k}(y_{k+1}\mid h_k,u_k)$，观测到 $y_{k+1}$ 后定义负对数似然 $\ell_k=-\log p_{\theta_k}(y_{k+1}\mid h_k,u_k)$。若数据身份和传感质量通过检查，可执行受约束更新
\[
\theta_{k+1}=\operatorname{Proj}_{\Theta_{\mathrm{safe}}}
\left(\theta_k-\alpha_k\nabla_{\theta}\ell_k\right).
\]
这是一个模型假设下的学习规则；$\ell_k$ 下降只表明模型对所见数据拟合改善，不推出任务必然成功，更不推出内部模型与环境整体同构。若残差来自目标改变、观测故障、接口延迟或其他行动者，直接更新 $\theta$ 反而会污染结构记忆。系统必须先做来源判别，再决定写入 $h(t)$ 的临时修正、写入 $E$ 的带证据事件、固化到 $\Theta$ 的长期约束，或卸载到 Me 的外部记录。

原稿提出“共形共生”作为演化终局。作者保留其任务相关尺度对齐的直觉，但拒绝全局终局与零摩擦承诺。定义任务特征 $z=q(s,e)$、内部预测 $\hat z=r(x,\theta)$、约束集合 $\mathcal C$ 与作用域 $D$。在时间窗 $[k,k+H]$ 内，如果预测误差、控制代价和风险同时受限，
\[
\sup_{j\in[k,k+H]}d(z_j,\hat z_j)\leq\varepsilon,
\qquad C^{\mathrm{act}}_{k:k+H}\leq B,
\qquad P(\mathcal C\text{ violated})\leq\delta,
\]
则称系统在 $(D,H,\varepsilon,B,\delta)$ 上达到任务相关共生。该定义允许局部尺度变换、注意增益和抽象压缩，却不要求无关背景与内部表示一一对应。窗口外、目标改变后或扰动超过设计域时，对齐可能立即失效；心流体验也不能替代误差、代价与风险的测量。

这种双向重构还解释了工具与外部记忆的地位。智能体若把重复计算外化为脚本，把稳定事实外化为数据库，把空间约束外化为建筑与标识，就改变了未来状态转移和观测函数。Offloading 因而不仅释放当前 $h(t)$，还可能降低以后任务对 $\Theta$ 内部复杂度的要求；Boundary 决定外化内容的访问性、稳定性、延迟与隐私风险。外化并非总是有利：工具失效、数据漂移或组织依赖会把代价推迟到未来。评估时必须把维护、召回、验证和失效恢复计入长期预算。环境既是被控制对象，也是记忆与计算的承载者，但这种具身扩展仍服从工作记忆容量与内外修改代价的约束。

最后，双向重构不意味着征服环境，也不意味着被动顺从。可靠系统在三类响应之间选择：当模型错误而环境约束稳定时更新内部；当环境可塑且改变具有长期收益时实施外部工程；当双方代价过高或风险不可接受时修改目标、缩小作用域或退出。组织智能还要考虑其他主体的权利与竞争目标，不能把他者简单当作阻抗。HSF-HD 在这里提供的是一般状态动力学语言：目标、状态、结构、边界和反馈如何共同变化；AgentOS 则可把它具体化为 TCE 调节、SPC 读出、CCM 负载管理、Boundary 授权和 Me 外部记忆。二者的关系是理论层与工程实例，而不是同义替换。至此，交互总线的判据不再是漂亮的几何比喻，而是动作可执行、传播可追踪、后果可观测、误差可归因、结构更新可审计。
